\section*{Bloque D. Ampliación: Teoría y Aplicaciones Prácticas (70\% / 30\%)}

\noindent\textit{Preguntas 91 a 180. Ampliación del banco oficial con 63 preguntas teóricas y conceptuales (91 a 153) y 27 aplicaciones cuantitativas prácticas (154 a 180).}

\subsection*{Parte I. Preguntas teóricas y conceptuales (91 a 153)}

\q{Una empresa que es tomadora de precios en el mercado laboral pero monopolista en el mercado del producto final contrata trabajo hasta el punto en que el salario $w$ iguala a:}
\begin{opts}
    \item El ingreso del producto marginal ($IPMg_L$), que es menor que el valor del producto marginal.
    \item El valor del producto marginal ($VPMg_L$), que coincide exactamente con el precio del bien.
    \item El producto medio del trabajo multiplicado por el costo marginal de producción de corto plazo.
    \item El costo medio variable del trabajo multiplicado por la elasticidad precio de la demanda final.
\end{opts}

\q{Se denomina «explotación monopólica» del factor trabajo (diferente de la explotación monopsonística) a la brecha positiva que surge cuando una empresa:}
\begin{opts}
    \item Opera en competencia perfecta en el producto y paga salarios iguales al valor del producto marginal.
    \item Tiene poder monopólico en el producto y paga al trabajo un salario inferior a su $VPMg_L$.
    \item Enfrenta una curva de oferta laboral creciente y contrata donde el salario iguala al producto medio.
    \item Aplica discriminación salarial de primer grado a trabajadores con idéntica productividad marginal.
\end{opts}

\q{¿Por qué la curva de demanda de trabajo de toda una industria competitiva es más inelástica (más empinada) que la suma horizontal de las curvas individuales de $VPMg_L$ de las empresas?}
\begin{opts}
    \item Porque la ley de rendimientos decrecientes desaparece al agregar la producción de varias empresas.
    \item Porque las empresas dejan de minimizar costos cuando todas enfrentan simultáneamente menor salario.
    \item Porque el menor salario expande la oferta del bien y cae $P$, desplazando hacia abajo $VPMg_L$.
    \item Porque el precio del capital aumenta automáticamente en la misma proporción que baja el salario.
\end{opts}

\q{Si trabajo y capital son insumos variables en el largo plazo y el trabajo es un factor normal, la curva de demanda de trabajo de largo plazo de una empresa competitiva es:}
\begin{opts}
    \item Más inelástica que la de corto plazo porque el capital fijo amplifica la respuesta del empleo.
    \item Idéntica a la de corto plazo siempre que la función de producción tenga retornos constantes.
    \item De pendiente positiva cuando la elasticidad de sustitución entre capital y trabajo supera uno.
    \item Más elástica que la de corto plazo al operar tanto el efecto sustitución como el de escala.
\end{opts}

\q{Ante una reducción en el salario $w$, ¿cómo actúan el efecto sustitución y el efecto escala (o producto) sobre la cantidad demandada de trabajo si el trabajo es un insumo normal?}
\begin{opts}
    \item Ambos efectos actúan en la misma dirección, incrementando la demanda de trabajo.
    \item El efecto sustitución aumenta el empleo mientras el efecto escala lo reduce necesariamente.
    \item El efecto sustitución es nulo en el largo plazo y solo opera el efecto escala sobre el producto.
    \item Ambos efectos se cancelan exactamente cuando la función de producción es homogénea de grado uno.
\end{opts}

\q{Si el trabajo fuera un insumo inferior (su demanda condicionada cae al aumentar el nivel de producto $Q$), ¿podría la curva de demanda no condicionada de trabajo de una firma competitiva tener pendiente positiva?}
\begin{opts}
    \item Sí, siempre que el efecto escala negativo supere en magnitud al efecto sustitución positivo.
    \item No, porque la convexidad de $\pi^*(p,w,r)$ garantiza $\partial L^*/\partial w \le 0$.
    \item Sí, pero únicamente cuando la tecnología de producción exhibe proporciones fijas de Leontief.
    \item No, porque los insumos inferiores solo existen en mercados con poder de monopolio en el bien.
\end{opts}

\q{De acuerdo con las leyes de Hicks-Marshall de la demanda derivada, manteniendo todo lo demás constante, la demanda de trabajo de una industria será más elástica cuanto:}
\begin{opts}
    \item Menor sea la elasticidad precio de la demanda del producto final que fabrica la industria.
    \item Menor sea la elasticidad de oferta de los demás factores productivos complementarios o sustitutos.
    \item Mayor sea la elasticidad de sustitución entre el trabajo y otros factores en la producción.
    \item Más rígida sea la tecnología de producción en proporciones fijas entre capital y trabajo.
\end{opts}

\q{La cuarta ley de Hicks-Marshall («la importancia de ser poco importante») indica que una mayor participación del trabajo en el costo total ($s_L$) eleva la elasticidad de la demanda de trabajo $|\varepsilon_{L,w}|$ siempre que:}
\begin{opts}
    \item La elasticidad de sustitución $\sigma$ sea estrictamente mayor que la elasticidad del producto $|\varepsilon_Q|$.
    \item La demanda del bien final sea completamente inelástica ($\varepsilon_Q = 0$) en el corto plazo.
    \item Los factores capital y trabajo se empleen en proporciones estrictamente variables con $\sigma = \infty$.
    \item La elasticidad precio del bien final $|\varepsilon_Q|$ supere a la elasticidad de sustitución $\sigma$.
\end{opts}

\q{Cuando sube el precio del capital $r$, la demanda no condicionada de trabajo de una empresa competitiva puede aumentar o disminuir dependiendo de la relación entre sustitución y escala. Se dice que trabajo y capital son «complementos brutos» si:}
\begin{opts}
    \item El efecto escala negativo supera al efecto sustitución positivo, reduciendo la demanda de trabajo.
    \item El efecto sustitución positivo supera al efecto escala negativo, aumentando la demanda de trabajo.
    \item La elasticidad de sustitución técnica a lo largo de la isocuanta es estrictamente negativa.
    \item El producto marginal del trabajo disminuye cuando aumenta el precio del producto final.
\end{opts}

\q{Según el lema de Shephard aplicado a la función de costos mínimos $C(w, r, Q)$ de una empresa, la demanda condicionada de trabajo $L^c(w, r, Q)$ se obtiene mediante:}
\begin{opts}
    \item La derivada parcial de la función de beneficios máximos respecto al precio del producto final.
    \item La derivada parcial de la función de costos mínimos respecto al salario $w$.
    \item El cociente entre el costo marginal de producción y la productividad media del capital.
    \item La elasticidad del costo total respecto al nivel de producción $Q$ multiplicada por el salario.
\end{opts}

\q{Para una empresa competitiva que maximiza beneficios con dos insumos ($L$ y $K$), la simetría de las derivadas segundas de la función de beneficios $\pi^*(p, w, r)$ implica que:}
\begin{opts}
    \item El efecto del salario sobre la demanda de trabajo es igual al efecto del precio de $K$ sobre $K$.
    \item Las elasticidades precio cruzadas entre capital y trabajo son siempre idénticas entre sí.
    \item Los efectos cruzados son simétricos: $\partial L^*/\partial r = \partial K^*/\partial w$.
    \item Un aumento simultáneo y proporcional de todos los precios duplica la demanda de trabajo.
\end{opts}

\q{Si el precio del bien final $P$, el salario $w$ y la tasa de alquiler del capital $r$ se duplican simultáneamente en un mercado competitivo, la demanda no condicionada de trabajo $L^*(P, w, r)$:}
\begin{opts}
    \item Se duplica exactamente por ser homogénea de grado uno en todos los precios.
    \item Se reduce a la mitad debido al encarecimiento proporcional de los costos variables.
    \item Aumenta solo si la función de producción exhibe rendimientos crecientes a escala.
    \item Permanece inalterada por ser homogénea de grado cero en el vector de precios.
\end{opts}

\q{En el largo plazo, una empresa competitiva que contrata trabajo y capital en mercados competitivos satisface en el óptimo interior la condición:}
\begin{opts}
    \item $\frac{PMg_L}{w} = \frac{PMg_K}{r} = \frac{1}{P}$, igualando el $PMg$ por sol gastado.
    \item $PMg_L = PMg_K = P$, independientemente de los precios relativos de ambos factores.
    \item $\frac{PMg_L}{r} = \frac{PMg_K}{w} = P$, maximizando la productividad media del capital invertido.
    \item $w \cdot PMg_L = r \cdot PMg_K = CMg$, minimizando únicamente el costo fijo medio.
\end{opts}

\q{En el modelo neoclásico de oferta de trabajo individual con dotación de tiempo $T$, si el ocio es un bien normal, la ecuación de Slutsky para las horas trabajadas $h$ frente a un aumento del salario $w$ muestra que:}
\begin{opts}
    \item Tanto el efecto sustitución como el efecto ingreso impulsan al trabajador a ofrecer más horas.
    \item El efecto sustitución eleva las horas trabajadas y el efecto ingreso las reduce.
    \item El efecto sustitución reduce las horas trabajadas y el efecto ingreso las incrementa.
    \item Ambos efectos son nulos siempre que el individuo carezca de ingresos no laborales.
\end{opts}

\q{La curva de oferta de trabajo individual se vuelve de pendiente negativa (doblada hacia atrás o backward-bending) para niveles salariales elevados cuando:}
\begin{opts}
    \item El ocio se transforma en un bien inferior y su consumo cae al crecer el ingreso.
    \item El salario de mercado cae por debajo del salario de reserva del trabajador representativo.
    \item El efecto ingreso negativo sobre las horas supera en magnitud al efecto sustitución.
    \item El impuesto proporcional sobre la renta laboral supera la tasa marginal de sustitución.
\end{opts}

\q{En el modelo ocio-consumo con ingreso no laboral $M > 0$, el salario de reserva $w_R$ de un individuo se define formalmente como:}
\begin{opts}
    \item El salario máximo que una empresa monopsonística está dispuesta a pagar por la primera hora.
    \item El cociente entre el ingreso no laboral $M$ y la dotación total de tiempo disponible $T$.
    \item El salario mínimo legal fijado por el gobierno para evitar el desempleo estructural.
    \item La tasa marginal de sustitución entre ocio y consumo evaluada en cero horas trabajadas.
\end{opts}

\q{Si el ocio es un bien normal y un trabajador recibe un aumento permanente en su ingreso no laboral $M$ (por ejemplo, una transferencia o herencia), ¿qué ocurre con su salario de reserva y sus horas trabajadas?}
\begin{opts}
    \item Su salario de reserva aumenta y, si continúa trabajando, reduce las horas ofrecidas.
    \item Su salario de reserva disminuye y aumenta las horas ofrecidas por efecto sustitución.
    \item Su salario de reserva permanece constante y sus horas trabajadas aumentan proporcionalmente.
    \item Su salario de reserva aumenta pero sus horas trabajadas no cambian al no variar el salario $w$.
\end{opts}

\q{Compare un impuesto proporcional sobre el salario (tasa $t$) con un impuesto de suma alzada (lump-sum) que recauda exactamente los mismos ingresos fiscales del mismo trabajador. Si el ocio es normal, bajo el impuesto proporcional el trabajador ofrecerá:}
\begin{opts}
    \item Más horas trabajadas porque el impuesto proporcional reduce el salario neto y genera efecto ingreso.
    \item Menos horas trabajadas porque el impuesto proporcional reduce el precio relativo del ocio.
    \item Exactamente las mismas horas trabajadas porque la recaudación y el ingreso disponible son iguales.
    \item Más horas trabajadas únicamente si la curva de oferta laboral tiene pendiente positiva.
\end{opts}

\q{Considere un trabajador que ya labora $h_0$ horas a un salario estándar $w$. Si la empresa ofrece una prima por horas extras $w_e > w$ únicamente para horas por encima de $h_0$ (en lugar de subir el salario de todas las horas a $w_e$), el trabajador:}
\begin{opts}
    \item Reducirá sus horas trabajadas si el ocio es un bien normal debido al gran efecto ingreso.
    \item Trabajará exactamente las mismas $h_0$ horas porque el salario marginal no cambia en el margen.
    \item Aumentará sus horas trabajadas porque en $h_0$ solo opera el efecto sustitución.
    \item Abandonará el mercado laboral al elevarse su salario de reserva por encima de $w_e$.
\end{opts}

\q{Si las preferencias de un individuo entre consumo $C$ y ocio $l$ son cuasilineales de la forma $U(C, l) = C + v(l)$ (con $v' > 0, v'' < 0$), entonces en una solución interior la curva de oferta de trabajo:}
\begin{opts}
    \item Se dobla hacia atrás necesariamente para cualquier nivel salarial positivo.
    \item Es completamente vertical porque los efectos ingreso y sustitución se cancelan.
    \item Depende negativamente del ingreso no laboral $M$ en todo el rango de ingresos.
    \item Tiene pendiente positiva porque el efecto ingreso sobre el ocio es cero.
\end{opts}

\q{Para un individuo con preferencias Cobb-Douglas $U(C, l) = C^\alpha l^{1-\alpha}$ y cero ingreso no laboral ($M = 0$), ¿cómo responde su oferta de horas trabajadas $h^*$ ante cambios en el salario $w$?}
\begin{opts}
    \item Permanece constante en $h^* = \alpha T$ porque los efectos ingreso y sustitución se compensan.
    \item Aumenta proporcionalmente con el salario $w$ para cualquier valor de $\alpha \in (0,1)$.
    \item Disminuye cuando aumenta el salario $w$ porque el efecto ingreso siempre domina al sustitución.
    \item Es igual a cero siempre que el salario sea inferior a la unidad monetaria.
\end{opts}

\q{Incluso si la oferta de trabajo de algunos trabajadores individuales se dobla hacia atrás a salarios altos, la oferta de trabajo agregada de una ocupación específica suele mantener pendiente positiva porque:}
\begin{opts}
    \item El ocio es un bien inferior para la mayoría de los trabajadores en economías modernas.
    \item Salarios más altos atraen nuevos participantes y trabajadores desde otras ocupaciones.
    \item La ecuación de Slutsky carece de validez cuando se agregan múltiples individuos.
    \item Los impuestos proporcionales eliminan el efecto ingreso a nivel agregado del mercado.
\end{opts}

\q{En el modelo intertemporal de oferta laboral a lo largo del ciclo de vida, la elasticidad de Frisch mide la respuesta de las horas trabajadas ante:}
\begin{opts}
    \item Un cambio permanente e inesperado en el salario real que altera la riqueza vitalicia.
    \item Un aumento en las transferencias no laborales manteniendo fija la tasa de interés.
    \item Un cambio salarial temporal manteniendo constante la utilidad marginal de la riqueza.
    \item Una variación en el salario mínimo legal en presencia de un empleador monopsonista.
\end{opts}

\q{En un programa de transferencia asistencial con garantía de ingreso mínimo que reduce el subsidio en $S/\ 1$ por cada $S/\ 1$ ganado trabajando (tasa de retiro del 100\%), el efecto teórico sobre la oferta laboral de los beneficiarios es:}
\begin{opts}
    \item Aumentar las horas trabajadas al incrementarse la productividad marginal del trabajador.
    \item Mantener inalteradas las horas trabajadas porque no cambia el salario bruto de mercado.
    \item Aumentar la participación laboral únicamente de quienes tienen salarios de reserva altos.
    \item Desincentivar el trabajo al imponer un salario neto marginal efectivo nulo.
\end{opts}

\q{En un monopsonio laboral con curva de oferta creciente $w(L)$ y sin discriminación salarial, el gasto marginal del trabajo $GMg_L$ supera al salario $w(L)$ porque:}
\begin{opts}
    \item Contratar un trabajador adicional exige subir el salario a todos los previos.
    \item La productividad marginal del trabajo crece a tasa acelerada conforme aumenta el empleo.
    \item El monopsonista enfrenta una demanda del producto final perfectamente inelástica.
    \item Los costos fijos de contratación se dividen entre un número cada vez menor de empleados.
\end{opts}

\q{Para un monopsonista competitivo en el mercado del producto que enfrenta una oferta laboral con elasticidad $\eta_L > 0$, la tasa de explotación monopsonística de Pigou-Robinson $\frac{VPMg_L - w}{w}$ en equilibrio equivale a:}
\begin{opts}
    \item El valor de la elasticidad de oferta laboral $\eta_L$ multiplicado por el salario.
    \item El inverso de la elasticidad de la oferta de trabajo ($1 / \eta_L$).
    \item La elasticidad precio de la demanda del producto final dividida por el salario.
    \item La diferencia entre el producto medio del trabajo y su costo fijo medio.
\end{opts}

\q{¿Bajo qué condición la imposición de un salario mínimo legal $w_{\min}$ en un mercado laboral monopsonístico incrementa simultáneamente el salario y el nivel de empleo respecto al equilibrio sin regulación?}
\begin{opts}
    \item Cuando $w_{\min}$ se fija por debajo del salario de equilibrio de monopsonio $w_m$.
    \item Cuando $w_{\min}$ se fija por encima del $VPMg_L$ correspondiente al nivel de empleo monopsonístico.
    \item Cuando $w_{\min}$ se ubica estrictamente entre $w_m$ y el $VPMg_L$ evaluado en $L_m$.
    \item Cuando la curva de oferta laboral que enfrenta el monopsonista es perfectamente elástica.
\end{opts}

\q{En un mercado monopsonístico, ¿qué nivel de salario mínimo $w_{\min}$ maximiza el nivel de empleo y elimina por completo la pérdida irrecuperable de eficiencia?}
\begin{opts}
    \item El salario que iguala el gasto marginal del factor con el producto medio del trabajo.
    \item El salario equivalente al valor del producto marginal evaluado en el empleo de monopsonio.
    \item Un salario mínimo igual al costo de oportunidad del primer trabajador contratado.
    \item El salario competitivo donde se intersecan las curvas de $VPMg_L$ y oferta laboral $w(L)$.
\end{opts}

\q{Un monopsonista que puede pagar salarios diferenciados a dos grupos de trabajadores igualmente productivos (grupos 1 y 2) con elasticidades de oferta laboral $\eta_1 < \eta_2$ fijará en equilibrio:}
\begin{opts}
    \item Un salario menor para el grupo 1 porque su oferta laboral es relativamente menos elástica.
    \item Un salario mayor para el grupo 1 porque su gasto marginal laboral es más bajo en el margen.
    \item El mismo salario para ambos grupos al tener exactamente la misma productividad marginal.
    \item Un empleo nulo en el grupo 2 para concentrar toda la contratación en el grupo inelástico.
\end{opts}

\q{Si un monopsonista logra practicar discriminación salarial de primer grado (pagando a cada trabajador exactamente su salario de reserva a lo largo de la curva de oferta), el resultado en términos de eficiencia asignativa es:}
\begin{opts}
    \item Una reducción del empleo aún mayor que en el monopsonio con salario uniforme.
    \item El mismo nivel de empleo que en competencia perfecta, eliminando la pérdida de eficiencia.
    \item Un aumento de la renta económica de los trabajadores a costa de los beneficios de la firma.
    \item Una brecha salarial positiva entre el último trabajador contratado y su $VPMg_L$.
\end{opts}

\q{En un monopolio bilateral (un sindicato monopolista que negocia con un único empleador monopsonista), ¿cómo se determinan el salario y el nivel de empleo de equilibrio?}
\begin{opts}
    \item El salario coincide necesariamente con el equilibrio competitivo sin importar la negociación.
    \item El empleo resulta siempre superior al nivel de competencia perfecta por presión sindical.
    \item El salario queda indeterminado dentro de un rango y depende del poder negociador.
    \item El sindicato impone siempre un salario igual al producto medio del trabajo en el corto plazo.
\end{opts}

\q{Al comparar el modelo sindical de «derecho a gestionar» (right-to-manage, donde se negocia solo el salario y la firma elige el empleo sobre su curva de demanda) con el modelo de «contratos eficientes» (se negocian salario y empleo conjuntamente):}
\begin{opts}
    \item El modelo de derecho a gestionar siempre se ubica sobre la curva de contrato Pareto-eficiente.
    \item Los contratos eficientes generan menor empleo y menor salario que el derecho a gestionar.
    \item Ambos modelos producen exactamente el mismo nivel de empleo cuando el sindicato es adverso al riesgo.
    \item Los contratos eficientes logran puntos sobre la curva de contrato, fuera de la demanda.
\end{opts}

\q{Si una empresa es simultáneamente monopolista en el mercado del producto y monopsonista en el mercado de trabajo, su condición de contratación óptima es:}
\begin{opts}
    \item $IPMg_L = GMg_L$, donde $IPMg_L < VPMg_L$ y $GMg_L > w$, reduciendo doblemente el empleo.
    \item $VPMg_L = w$, porque ambos poderes de mercado se compensan mutuamente en el óptimo.
    \item $IPMg_L = w$, ya que el poder monopólico anula la pendiente positiva de la oferta laboral.
    \item $VPMg_L = GMg_L$, generando una pérdida de eficiencia idéntica a la del monopsonio puro.
\end{opts}

\q{De manera análoga a la ausencia de curva de oferta en un monopolio, ¿por qué un monopsonista carece de una curva de demanda de trabajo bien definida?}
\begin{opts}
    \item Porque su función de producción presenta rendimientos crecientes a escala en el corto plazo.
    \item Porque el empleo depende de la forma y elasticidad de la oferta, no solo de $w$.
    \item Porque el salario pagado en un monopsonio es siempre estrictamente superior al $VPMg_L$.
    \item Porque no existe una relación funcional entre el nivel de empleo y la producción física.
\end{opts}

\q{En la economía laboral moderna (modelo de Manning / Burdett-Mortensen), ¿por qué muchas empresas poseen cierto poder monopsonístico incluso en ciudades con numerosos empleadores?}
\begin{opts}
    \item Porque la ley prohíbe a los trabajadores cambiar de empleador sin autorización judicial previa.
    \item Porque todos los trabajadores poseen exactamente las mismas preferencias por ubicación y horarios.
    \item Porque existen fricciones de búsqueda, costos de movilidad y diferenciación de puestos de trabajo.
    \item Porque las empresas pequeñas tienen curvas de oferta laboral perfectamente horizontales.
\end{opts}

\q{En un mercado laboral perfectamente competitivo con oferta de pendiente positiva y demanda de pendiente negativa, un salario mínimo vinculante $w_{\min} > w^*$ genera:}
\begin{opts}
    \item Un aumento simultáneo del empleo contratado y de la renta económica de todos los trabajadores.
    \item Un exceso de demanda de trabajo (escasez de mano de obra) y eliminación del desempleo.
    \item Ningún efecto sobre el empleo siempre que la elasticidad de demanda laboral sea finita.
    \item Menor empleo contratado, desempleo involuntario y pérdida irrecuperable de eficiencia.
\end{opts}

\q{Un sindicato monopolista que enfrenta la curva de demanda laboral competitiva de una industria y busca maximizar la masa salarial total ($w \cdot L$) elegirá el salario donde:}
\begin{opts}
    \item La elasticidad precio de la demanda de trabajo sea igual a $-1$ (ingreso marginal laboral nulo).
    \item El salario iguale exactamente al salario de reserva del trabajador menos productivo.
    \item La demanda de trabajo se encuentre en su tramo estrictamente inelástico ($|\varepsilon_{L,w}| < 1$).
    \item El nivel de empleo sea idéntico al del equilibrio competitivo sin intervención sindical.
\end{opts}

\q{La renta económica pura de un factor productivo se define como la diferencia entre el pago total que recibe el factor y:}
\begin{opts}
    \item El costo fijo promedio de la empresa que contrata dicho factor en el largo plazo.
    \item El pago mínimo necesario (pago de transferencia) para mantenerlo en su uso actual.
    \item El ingreso total generado por la venta del producto final en el mercado competitivo.
    \item El valor presente de los impuestos sobre la renta pagados por el propietario del factor.
\end{opts}

\q{Si la oferta de un factor productivo (como la tierra urbana no reproducible) es perfectamente inelástica, el 100\% de su remuneración constituye renta económica pura. Un impuesto sobre dicha renta:}
\begin{opts}
    \item Se traslada íntegramente a los consumidores finales elevando el precio de los bienes.
    \item Reduce drásticamente la cantidad ofrecida de tierra y genera una alta pérdida de eficiencia.
    \item Recae totalmente sobre los propietarios y no genera pérdida de eficiencia.
    \item Provoca la salida inmediata de la tierra hacia usos alternativos más rentables.
\end{opts}

\q{Según Alfred Marshall, la «cuasi-renta» generada por un activo de capital fijo en el corto plazo es igual a:}
\begin{opts}
    \item El beneficio económico neto después de deducir todos los costos fijos y variables.
    \item El costo fijo total dividido entre el número de unidades producidas en el periodo.
    \item La diferencia entre el precio de mercado y el costo marginal multiplicada por la tasa de interés.
    \item El exceso del ingreso total sobre los costos variables totales ($IT - CVT$).
\end{opts}

\q{En el corto plazo, una empresa competitiva con costos fijos hundidos decide continuar operando (en lugar de cerrar temporalmente) siempre que su cuasi-renta sea:}
\begin{opts}
    \item Mayor o igual a cero ($IT \ge CVT$), aun con beneficios económicos negativos.
    \item Estrictamente superior a sus costos fijos totales ($IT - CVT > CF$).
    \item Exactamente igual a cero, garantizando beneficios económicos nulos en el corto plazo.
    \item Negativa pero menor en valor absoluto que los costos fijos hundidos de la planta.
\end{opts}

\q{El teorema de Euler aplicado a la teoría de la distribución marginalista demuestra que si cada factor recibe su productividad marginal real ($PMg_i$), el producto total $Q$ se agota exactamente ($Q = L\cdot PMg_L + K\cdot PMg_K$) si y solo si:}
\begin{opts}
    \item La función de producción exhibe rendimientos crecientes a escala en todo su dominio.
    \item La función de producción tiene rendimientos constantes a escala (grado uno).
    \item La elasticidad de sustitución entre capital y trabajo es estrictamente igual a cero.
    \item Los mercados de bienes operan bajo competencia monopolística con exceso de capacidad.
\end{opts}

\q{Si una empresa opera con una función de producción homogénea de grado $k > 1$ (rendimientos crecientes a escala) y remunerara a cada factor exactamente con su valor del producto marginal ($VPMg_i$), entonces:}
\begin{opts}
    \item Obtendría beneficios económicos positivos iguales a $(k - 1)PQ$ en el equilibrio.
    \item El ingreso total se agotaría exactamente dejando beneficios económicos iguales a cero.
    \item Los pagos totales a factores superarían al ingreso total, generando pérdidas.
    \item La participación del trabajo en el ingreso total sería necesariamente igual a cero.
\end{opts}

\q{En una economía competitiva con función de producción agregada Cobb-Douglas $Y = A L^\alpha K^{1-\alpha}$, si se acumula capital aumentando la razón $K/L$ mientras la tecnología es constante, la participación del trabajo en el ingreso total ($wL / Y$):}
\begin{opts}
    \item Aumenta progresivamente porque el trabajo se vuelve relativamente más escaso que el capital.
    \item Disminuye progresivamente porque el capital sustituye al trabajo en la producción agregada.
    \item Converge a cero en el largo plazo conforme la razón capital-trabajo tiende a infinito.
    \item Permanece constante e igual a $\alpha$ porque la elasticidad de sustitución es $\sigma = 1$.
\end{opts}

\q{Considere una economía competitiva con función de producción CES cuya elasticidad de sustitución entre capital y trabajo satisface $\sigma > 1$ (factores fácilmente sustituibles). Si la razón capital-trabajo $K/L$ aumenta:}
\begin{opts}
    \item La participación del capital en el ingreso aumenta y la participación del trabajo disminuye.
    \item La participación del trabajo en el ingreso aumenta y la participación del capital disminuye.
    \item Las participaciones del capital y del trabajo permanecen estrictamente inalteradas.
    \item El salario real disminuye en términos absolutos aunque aumente la productividad del trabajo.
\end{opts}

\q{En una función de producción CES competitiva con elasticidad de sustitución $\sigma < 1$ (capital y trabajo relativamente complementarios), un incremento en la razón capital-trabajo $K/L$ ocasiona que:}
\begin{opts}
    \item La participación del capital en el ingreso nacional aumente a costa de la participación salarial.
    \item La participación del trabajo en el ingreso nacional aumente y la del capital disminuya.
    \item El producto marginal del trabajo disminuya por efecto de los rendimientos decrecientes.
    \item Los precios relativos de los factores permanezcan fijos independientemente de $K/L$.
\end{opts}

\q{Si un mercado competitivo enfrenta una oferta de trabajo perfectamente elástica (horizontal) al salario de mercado $w^*$, ¿a cuánto asciende la renta económica percibida por los trabajadores?}
\begin{opts}
    \item A la totalidad de la masa salarial $w^* L^*$ pagada por las empresas contratantes.
    \item A la mitad del excedente total generado en el mercado de factores productivo.
    \item A cero, pues el salario coincide con el costo de oportunidad del trabajador.
    \item Al beneficio económico extraordinario obtenido por las empresas en el corto plazo.
\end{opts}

\q{Para un profesional altamente especializado (por ejemplo, un ingeniero sanitario experto en plantas de tratamiento), la proporción de su salario que constituye renta económica suele ser:}
\begin{opts}
    \item Nula en el corto plazo y equivalente al 100\% en el largo plazo por la libre entrada.
    \item Idéntica en el corto y en el largo plazo independientemente de la movilidad sectorial.
    \item Mayor en el largo plazo porque la curva de oferta laboral especializada se vuelve vertical.
    \item Mayor en el corto plazo que en el largo plazo por la menor elasticidad de oferta.
\end{opts}

\q{Según Gary Becker, en mercados laborales competitivos, ¿quién financia normalmente la capacitación general (útil por igual en cualquier empresa) frente a la capacitación específica (útil solo en la empresa actual)?}
\begin{opts}
    \item El trabajador financia la general (vía menor salario inicial) y ambos comparten la específica.
    \item La empresa financia el 100\% de la general y el trabajador financia el 100\% de la específica.
    \item El Estado financia obligatoriamente ambas porque las empresas nunca comparten costos específicos.
    \item La empresa financia la general mediante bonos y recupera el costo pagando menos del $VPMg$ luego.
\end{opts}

\q{En el modelo de señalización del mercado laboral de Michael Spence, ¿qué condición técnica es indispensable para que exista un equilibrio separador donde la educación revele la productividad innata sin aumentarla?}
\begin{opts}
    \item Que la educación incremente el producto marginal de todos los trabajadores en igual proporción.
    \item Que el costo marginal educativo sea menor para el tipo de alta productividad.
    \item Que las empresas observen directamente la habilidad innata antes de fijar los salarios.
    \item Que el costo de estudiar sea idéntico para todos los tipos pero el salario sea regulado.
\end{opts}

\q{En un modelo puro de señalización de Spence (donde la educación no aumenta la productividad real sino solo actúa como señal), el equilibrio separador respecto a un equilibrio agrupador (pooling) puede resultar:}
\begin{opts}
    \item Siempre superior en sentido de Pareto porque todos los trabajadores ganan mayores salarios.
    \item Siempre libre de costos sociales porque la inversión educativa es una transferencia pura.
    \item Socialmente costoso (incluso Pareto-inferior) al gastar recursos en la señal.
    \item Imposible de sostener si existen más de dos empresas compitiendo por los trabajadores.
\end{opts}

\q{En el modelo de salarios de eficiencia por riesgo moral (Shapiro-Stiglitz), ¿por qué las empresas pagan voluntariamente salarios superiores al nivel que vacía el mercado, generando desempleo involuntario de equilibrio?}
\begin{opts}
    \item Porque los sindicatos obligan legalmente a fijar salarios por encima del $VPMg_L$.
    \item Porque la ley de rendimientos decrecientes impide contratar trabajadores a salarios más bajos.
    \item Porque las empresas desconocen el nivel de salario que equilibra la oferta y demanda agregada.
    \item Porque el salario alto y el desempleo elevan el costo del despido y evitan pereza.
\end{opts}

\q{Según la teoría de los diferenciales salariales compensatorios de Adam Smith y Sherwin Rosen, si dos empleos requieren igual calificación pero uno conlleva mayor riesgo laboral o peores condiciones ambientales:}
\begin{opts}
    \item El empleo más riesgoso ofrecerá un salario mayor para compensar la desutilidad del riesgo.
    \item Ambos empleos pagarán exactamente el mismo salario en un mercado competitivo de largo plazo.
    \item El empleo más seguro pagará un salario mayor porque atrae trabajadores más productivos.
    \item El diferencial salarial desaparecerá completamente cuando todos los trabajadores sean adversos al riesgo.
\end{opts}

\q{En la teoría de compensación basada en torneos (Lazear y Rosen), donde los premios salariales dependen del desempeño relativo (ranking) en lugar de la producción absoluta, el esquema es especialmente útil cuando:}
\begin{opts}
    \item La producción individual es medible con exactitud sin costo y no existen choques comunes.
    \item Medir el desempeño relativo es más barato que el absoluto y existen choques comunes de riesgo.
    \item Los trabajadores pueden sabotear fácilmente el desempeño de sus compañeros sin ser detectados.
    \item La empresa opera con un único trabajador y desea eliminar por completo el riesgo salarial.
\end{opts}

\q{Según el modelo de contratos de incentivos de Lazear (1979), ¿por qué las empresas pueden ofrecer perfiles salariales con pendiente más pronunciada que la productividad (pagando menos que el $VPMg$ a jóvenes y más a veteranos)?}
\begin{opts}
    \item Porque la productividad física de los trabajadores siempre crece indefinidamente hasta el retiro.
    \item Para incentivar la rotación temprana de personal y evitar costos de capacitación específica.
    \item Para actuar como fianza que desalienta el bajo esfuerzo y reduce la rotación voluntaria.
    \item Porque las leyes laborales prohíben pagar salarios proporcionales al producto marginal inicial.
\end{opts}

\q{La diferencia teórica central entre la discriminación laboral por gustos (Becker) y la discriminación estadística (Phelps-Arrow) radica en que la discriminación estadística:}
\begin{opts}
    \item Se origina en prejuicios irracionales del empleador y tiende a desaparecer bajo competencia perfecta.
    \item Implica siempre menores beneficios para las empresas que la practican frente a sus competidoras.
    \item Solo ocurre cuando existe un sindicato monopolista que controla la contratación del sector.
    \item Surge de información imperfecta usando promedios grupales y puede persistir en competencia.
\end{opts}

\q{Al estimar econométricamente el retorno privado de la educación mediante una ecuación de Mincer por MCO, el «sesgo de habilidad» (ability bias) ocurre porque:}
\begin{opts}
    \item La habilidad innata no observada se correlaciona positivamente con educación y salarios.
    \item Los años de escolaridad se miden siempre sin error en las encuestas de hogares.
    \item La experiencia laboral potencial disminuye automáticamente al aumentar la edad del individuo.
    \item El capital humano específico reduce la productividad marginal de la educación formal.
\end{opts}

\q{En la variante de selección adversa de los salarios de eficiencia (Weiss), una empresa que enfrenta exceso de postulantes se resiste a recortar el salario ofrecido porque:}
\begin{opts}
    \item El recorte salarial aumenta inmediatamente la productividad promedio de los postulantes.
    \item Bajar el salario empeora la calidad promedio al retirarse los mejores candidatos.
    \item La curva de oferta de trabajadores de alta calificación es completamente inelástica al salario.
    \item Los costos fijos de selección se reducen a cero cuando el salario supera el equilibrio.
\end{opts}

\q{En la teoría neoclásica de la demanda de capital (Jorgenson), una empresa competitiva que compra capital a precio $P_K$, enfrenta una tasa de interés real $r$ y una depreciación $\delta$ (sin impuestos ni ganancias de capital) contrata capital hasta que $VPMg_K$ iguala a:}
\begin{opts}
    \item $P_K \cdot r$, ignorando la depreciación porque es un gasto estrictamente contable.
    \item $P_K / (r + \delta)$, equivalente al valor presente perpetuo del flujo de capital.
    \item $P_K (r + \delta)$, sumando el costo financiero y la depreciación física.
    \item $P_K (r - \delta)$, descontando el desgaste físico de la tasa de interés de mercado.
\end{opts}

\q{Según la regla de Hotelling para un recurso natural agotable extraído en competencia perfecta con costo marginal constante $c$ y tasa de interés $r$, en equilibrio intertemporal:}
\begin{opts}
    \item El precio bruto de mercado permanece constante en el tiempo hasta el agotamiento físico.
    \item El costo marginal de extracción crece a una tasa exactamente igual a la tasa de interés $r$.
    \item La cantidad extraída aumenta exponencialmente a la tasa $r$ durante todo el horizonte.
    \item La renta de escasez ($P_t - c$) crece a la tasa de interés de mercado $r$.
\end{opts}

\q{En la gestión económica de recursos hídricos subterráneos (acuíferos sobreexplotados que abastecen a una EPS), el precio eficiente del agua debe incluir, además del costo marginal de bombeo y tratamiento:}
\begin{opts}
    \item El costo de uso intertemporal (renta de escasez) por menor disponibilidad futura.
    \item Únicamente el costo histórico amortizado de las tuberías secundarias de distribución.
    \item Un subsidio negativo igual a la depreciación contable de los medidores domiciliarios.
    \item El índice de Lerner correspondiente a un monopolio no regulado de bienes duraderos.
\end{opts}

\q{En una empresa prestadora de servicios de saneamiento (EPS) regulada por SUNASS, la existencia de sobrecostos laborales (rigideces de planilla o salarios por encima de niveles eficientes de mercado) impacta la regulación tarifaria porque:}
\begin{opts}
    \item SUNASS está obligada a trasladar automáticamente el 100\% de la planilla histórica a la tarifa.
    \item Al reconocerse solo el OPEX eficiente, los sobrecostos no cubiertos deterioran la caja de la EPS.
    \item Los salarios de la EPS se excluyen completamente del cálculo del flujo de caja del PMO.
    \item El efecto Averch-Johnson obliga a la EPS a contratar exclusivamente más personal en lugar de capital.
\end{opts}

\q{En la determinación de la tasa de descuento regulatoria (WACC) que remunera el factor capital en una empresa de saneamiento regulada por SUNASS, un aumento en el riesgo país o en el beta del sector:}
\begin{opts}
    \item Reduce el costo del capital propio y disminuye la tarifa media requerida en el quinquenio.
    \item No afecta el costo del capital porque las EPS son monopolios naturales sin riesgo sistemático.
    \item Eleva el costo del capital propio (CAPM), aumentando el WACC regulatorio.
    \item Obliga a sustituir el costo del capital propio por la tasa de depreciación física de las redes.
\end{opts}

\subsection*{Parte II. Aplicaciones prácticas y ejercicios cuantitativos (154 a 180)}

\q{Una empresa competitiva en los mercados de producto y trabajo tiene la función de producción de corto plazo $Q = 20L - 0.5L^2$. Si el precio del producto es $P = 10$ y el salario de mercado es $w = 40$, ¿cuál es el empleo óptimo $L^*$ y el beneficio bruto ($PQ - wL$)?}
\begin{opts}
    \item $L^* = 12$ y beneficio bruto de $1080$
    \item $L^* = 20$ y beneficio bruto de $1200$
    \item $L^* = 16$ y beneficio bruto de $1440$
    \item $L^* = 16$ y beneficio bruto de $1280$
\end{opts}

\q{Un monopolista en el mercado del producto (pero tomador de salarios en el mercado laboral) produce con $Q = 4L$ y enfrenta la demanda $P = 100 - Q$. Si el salario es $w = 80$, ¿cuáles son el empleo óptimo $L^*$ y el valor del producto marginal $VPMg_L$ en el equilibrio?}
\begin{opts}
    \item $L^* = 10$ y $VPMg_L = 240$
    \item $L^* = 20$ y $VPMg_L = 80$
    \item $L^* = 10$ y $VPMg_L = 80$
    \item $L^* = 15$ y $VPMg_L = 160$
\end{opts}

\q{Una empresa competitiva opera con la función de producción $Q = L^{0.5}K^{0.5}$ y un capital fijo de corto plazo $K = 16$. Si el precio del producto es $P = 20$ y el salario es $w = 10$, ¿cuántas unidades de trabajo $L^*$ contratará?}
\begin{opts}
    \item $L^* = 8$ unidades
    \item $L^* = 16$ unidades
    \item $L^* = 20$ unidades
    \item $L^* = 32$ unidades
\end{opts}

\q{Una empresa minimizadora de costos opera con la tecnología $Q = L^{0.4}K^{0.6}$ y enfrenta precios de factores $w = 30$ y $r = 20$. ¿Cuál es la razón óptima capital-trabajo $K/L$ en el largo plazo?}
\begin{opts}
    \item $K / L = 1.50$
    \item $K / L = 1.00$
    \item $K / L = 2.25$
    \item $K / L = 0.67$
\end{opts}

\q{Un trabajador con preferencias $U(C, l) = C^{0.6}l^{0.4}$ dispone de $T = 20$ horas diarias, carece de ingreso no laboral ($M = 0$) y enfrenta un salario de $w = 15$ por hora (con $P_C = 1$). ¿Cuántas horas diarias trabaja $h^*$ y cuál es su consumo $C^*$?}
\begin{opts}
    \item $h^* = 8$ horas y $C^* = 120$
    \item $h^* = 10$ horas y $C^* = 150$
    \item $h^* = 14$ horas y $C^* = 210$
    \item $h^* = 12$ horas y $C^* = 180$
\end{opts}

\q{Un individuo tiene utilidad $U(C, l) = C^{0.5}l^{0.5}$, dotación de tiempo $T = 16$ horas, salario $w = 10$ por hora e ingreso no laboral $M = 40$ (con $P_C = 1$). ¿Cuántas horas ofrece al mercado laboral?}
\begin{opts}
    \item $h^* = 6$ horas
    \item $h^* = 8$ horas
    \item $h^* = 10$ horas
    \item $h^* = 4$ horas
\end{opts}

\q{Un individuo con preferencias $U(C, l) = C \cdot l$ dispone de $T = 16$ horas diarias y recibe un ingreso no laboral de $M = 80$ diarios (con $P_C = 1$). ¿Cuál es su salario de reserva $w_R$?}
\begin{opts}
    \item $w_R = 10.00$ por hora
    \item $w_R = 5.00$ por hora
    \item $w_R = 8.00$ por hora
    \item $w_R = 4.00$ por hora
\end{opts}

\q{Un trabajador que labora $h = 40$ horas semanales exhibe un efecto total del salario $\frac{\partial h}{\partial w} = 0.40$ y un efecto del ingreso no laboral $\frac{\partial h}{\partial M} = -0.02$. Según la ecuación de Slutsky, ¿cuál es el efecto sustitución compensado $\left.\frac{\partial h}{\partial w}\right|_{\bar{U}}$?}
\begin{opts}
    \item $0.80$
    \item $0.40$
    \item $1.20$
    \item $0.60$
\end{opts}

\q{Un mercado de trabajo competitivo tiene demanda $L_d = 200 - 4w$ y oferta $L_s = -40 + 8w$. ¿Cuáles son el salario de equilibrio $w^*$, el empleo $L^*$ y la renta económica total de los trabajadores?}
\begin{opts}
    \item $w^* = 20$, $L^* = 120$, renta de $1200$
    \item $w^* = 25$, $L^* = 100$, renta de $1000$
    \item $w^* = 15$, $L^* = 140$, renta de $700$
    \item $w^* = 20$, $L^* = 120$, renta de $900$
\end{opts}

\q{Un empleador monopsonista enfrenta la oferta de trabajo $w(L) = 10 + L$ y tiene un valor del producto marginal $VPMg_L = 100 - L$. ¿Cuáles son el nivel de empleo $L_m^*$ y el salario $w_m^*$ de equilibrio?}
\begin{opts}
    \item $L_m^* = 30$ y $w_m^* = 40$
    \item $L_m^* = 45$ y $w_m^* = 55$
    \item $L_m^* = 30$ y $w_m^* = 70$
    \item $L_m^* = 40$ y $w_m^* = 50$
\end{opts}

\q{Un monopsonista enfrenta la oferta laboral $w(L) = 20 + L$ y presenta $VPMg_L = 140 - L$. ¿Cuál es la pérdida irrecuperable de eficiencia social (PBE) generada por el monopsonio respecto al equilibrio competitivo?}
\begin{opts}
    \item $800$ unidades monetarias
    \item $400$ unidades monetarias
    \item $200$ unidades monetarias
    \item $600$ unidades monetarias
\end{opts}

\q{Un monopsonista opera en un punto donde el valor del producto marginal del trabajo es $VPMg_L = 70$ y la elasticidad de la oferta laboral es $\eta_L = 2.5$. ¿Cuáles son el salario pagado $w^*$ y la tasa de explotación monopsonística $\frac{VPMg_L - w^*}{w^*}$?}
\begin{opts}
    \item $w^* = 42.0$ y tasa del $66.7\%$
    \item $w^* = 56.0$ y tasa del $25.0\%$
    \item $w^* = 50.0$ y tasa del $40.0\%$
    \item $w^* = 45.0$ y tasa del $55.6\%$
\end{opts}

\q{En el monopsonio con oferta laboral $w(L) = 10 + L$ y $VPMg_L = 100 - L$ (cuyo equilibrio libre es $L_m = 30, w_m = 40$), el gobierno fija un salario mínimo $w_{\min} = 50$. ¿Cuál es el nuevo nivel de empleo efectivo y su variación respecto a $L_m$?}
\begin{opts}
    \item $L = 50$ y el empleo aumenta en $20$
    \item $L = 30$ y el empleo permanece igual
    \item $L = 45$ y el empleo aumenta en $15$
    \item $L = 40$ y el empleo aumenta en $10$
\end{opts}

\q{Un monopsonista enfrenta $w(L) = 20 + L$ y $VPMg_L = 110 - L$. ¿Cuál es el salario mínimo máximo que puede imponer el gobierno sin que el nivel de empleo caiga por debajo del equilibrio de monopsonio no regulado?}
\begin{opts}
    \item $w_{\min} = 80$
    \item $w_{\min} = 65$
    \item $w_{\min} = 50$
    \item $w_{\min} = 90$
\end{opts}

\q{Una empresa es monopolista en el bien final ($P = 100 - Q$) y monopsonista en el trabajo ($w = 10 + 0.5L$), con función de producción $Q = L$. ¿Cuáles son el empleo $L^*$, el salario $w^*$ y el precio del producto $P^*$?}
\begin{opts}
    \item $L^* = 45$, $w^* = 32.5$, $P^* = 55$
    \item $L^* = 30$, $w^* = 25.0$, $P^* = 70$
    \item $L^* = 30$, $w^* = 40.0$, $P^* = 70$
    \item $L^* = 60$, $w^* = 40.0$, $P^* = 40$
\end{opts}

\q{Un monopsonista con $VPMg_L = 50$ constante contrata dos grupos separados de trabajadores con ofertas $w_1 = 10 + L_1$ y $w_2 = 20 + 0.5L_2$. Si aplica discriminación salarial, ¿qué salarios $(w_1^*, w_2^*)$ pagará a cada grupo?}
\begin{opts}
    \item $w_1^* = 35$ y $w_2^* = 30$
    \item $w_1^* = 40$ y $w_2^* = 40$
    \item $w_1^* = 30$ y $w_2^* = 35$
    \item $w_1^* = 20$ y $w_2^* = 30$
\end{opts}

\q{En un mercado laboral competitivo con curva de oferta $w = 20 + 0.5L$, el salario de equilibrio es $w^* = 50$. ¿A cuánto ascienden la renta económica y el pago de transferencia total, respectivamente?}
\begin{opts}
    \item $1200$ y $1800$
    \item $1500$ y $1500$
    \item $2100$ y $900$
    \item $900$ y $2100$
\end{opts}

\q{En el corto plazo, una empresa competitiva vende $Q = 100$ unidades a un precio $P = 15$, incurriendo en costos variables totales $CVT = 900$ y costos fijos $CF = 400$. ¿A cuánto ascienden la cuasi-renta y el beneficio económico?}
\begin{opts}
    \item Cuasi-renta de $600$ y beneficio de $200$
    \item Cuasi-renta de $200$ y beneficio de $600$
    \item Cuasi-renta de $1100$ y beneficio de $200$
    \item Cuasi-renta de $600$ y beneficio de $400$
\end{opts}

\q{Una empresa competitiva produce con $Q = 10L^{0.6}K^{0.4}$ empleando $L = 10$ y $K = 10$, y vende su producto a $P = 2$. Si cada factor recibe su $VPMg$, ¿cuánto suman la masa salarial ($wL$) y la remuneración al capital ($rK$)?}
\begin{opts}
    \item $wL = 100$ y $rK = 100$
    \item $wL = 120$ y $rK = 80$
    \item $wL = 80$ y $rK = 120$
    \item $wL = 150$ y $rK = 50$
\end{opts}

\q{Suponga una tecnología dada por $Q = L^{0.7}K^{0.5}$. Si hipotéticamente se remunerara tanto al trabajo como al capital con su respectivo valor del producto marginal ($VPMg_i$), ¿qué porcentaje del ingreso total $PQ$ representarían los pagos totales a los factores?}
\begin{opts}
    \item $100\%$ del ingreso total
    \item $80\%$ del ingreso total
    \item $120\%$ del ingreso total
    \item $140\%$ del ingreso total
\end{opts}

\q{En una economía competitiva descrita por la función de producción agregada $Y = A L^{0.60} K^{0.40}$ donde cada factor percibe su productividad marginal, ¿cuál es la razón entre la masa salarial y la remuneración del capital ($wL / rK$)?}
\begin{opts}
    \item $wL / rK = 0.67$
    \item $wL / rK = 1.25$
    \item $wL / rK = 2.00$
    \item $wL / rK = 1.50$
\end{opts}

\q{En una industria competitiva (con oferta de capital perfectamente elástica), la participación del trabajo en el costo total es $s_L = 0.40$, la elasticidad de sustitución es $\sigma = 0.50$ y la elasticidad precio de la demanda del bien es $|\varepsilon_Q| = 1.50$. Según Hicks-Marshall, ¿cuál es la elasticidad de la demanda de trabajo $|\varepsilon_{L,w}|$?}
\begin{opts}
    \item $|\varepsilon_{L,w}| = 0.90$
    \item $|\varepsilon_{L,w}| = 0.60$
    \item $|\varepsilon_{L,w}| = 1.10$
    \item $|\varepsilon_{L,w}| = 1.00$
\end{opts}

\q{Un profesional vive dos periodos ($t = 0, 1$) y descuenta a $r = 10\%$. Sin posgrado gana $S/\ 40\ 000$ en cada periodo. Si realiza una maestría en $t = 0$, paga $S/\ 10\ 000$ de matrícula y deja de trabajar ese periodo. ¿Cuál es el salario mínimo $X$ en $t = 1$ para que la maestría sea financieramente rentable?}
\begin{opts}
    \item $S/\ 90\ 000$
    \item $S/\ 95\ 000$
    \item $S/\ 85\ 000$
    \item $S/\ 55\ 000$
\end{opts}

\q{En un modelo de señalización de Spence, los trabajadores de alta productividad producen $a_H = 100$ (con costo educativo $10e$) y los de baja productividad producen $a_L = 50$ (con costo educativo $25e$). ¿Qué intervalo de años de educación $e^*$ sostiene un equilibrio separador?}
\begin{opts}
    \item $1 \le e^* \le 4$
    \item $2 \le e^* \le 4$
    \item $2 \le e^* \le 5$
    \item $3 \le e^* \le 5$
\end{opts}

\q{Una empresa adquiere maquinaria a un precio unitario $P_K = 200$, enfrenta una tasa de interés real de $r = 6\%$ anual y una tasa de depreciación física de $\delta = 9\%$ anual (sin impuestos ni inflación). ¿Cuál es el costo anual de uso del capital $c_K$?}
\begin{opts}
    \item $c_K = 12.00$
    \item $c_K = 18.00$
    \item $c_K = 24.00$
    \item $c_K = 30.00$
\end{opts}

\q{Un recurso natural agotable se extrae en competencia perfecta con costo marginal constante. Si el precio neto (renta de escasez $P_0 - c$) hoy es de $40.00$ y la tasa de interés de mercado es $r = 5\%$ anual, ¿cuál será el precio neto esperado dentro de $2$ años según la regla de Hotelling?}
\begin{opts}
    \item $44.10$ unidades monetarias
    \item $42.00$ unidades monetarias
    \item $44.00$ unidades monetarias
    \item $48.40$ unidades monetarias
\end{opts}

\q{Una planta de tratamiento de agua de una EPS produce $Q = 100\sqrt{L}$ metros cúbicos diarios con $L$ operadores técnicos. Si la tarifa regulada es $P = S/\ 4/\text{m}^3$ y el salario diario por operador es $w = S/\ 100$, ¿cuál es el número óptimo de operadores $L^*$ y la producción diaria $Q^*$?}
\begin{opts}
    \item $L^* = 2$ operadores y $Q^* = 141\ \text{m}^3$
    \item $L^* = 4$ operadores y $Q^* = 200\ \text{m}^3$
    \item $L^* = 8$ operadores y $Q^* = 283\ \text{m}^3$
    \item $L^* = 16$ operadores y $Q^* = 400\ \text{m}^3$
\end{opts}
